\documentclass[UTF8,11pt]{ctexart}
\usepackage[a4paper,margin=24mm]{geometry}
\usepackage{amsmath,amssymb,amsthm,mathtools}
\usepackage[all]{xy}
\usepackage{xurl,hyperref,enumitem}
\hypersetup{hidelinks,breaklinks=true}
\setlength{\emergencystretch}{3em}
\newtheorem*{theorem}{Theorem}
\newtheorem*{lemma}{Lemma}
\newtheorem*{proposition}{Proposition}
\newtheorem*{corollary}{Corollary}
\theoremstyle{definition}
\newtheorem*{definition}{Definition}
\newtheorem*{remark}{Remark}
\DeclareMathOperator{\Hom}{Hom}
\DeclareMathOperator{\Ext}{Ext}
\DeclareMathOperator{\Tor}{Tor}
\DeclareMathOperator{\Spec}{Spec}
\DeclareMathOperator{\supp}{supp}
\DeclareMathOperator{\im}{im}
\DeclareMathOperator{\id}{id}
\DeclareMathOperator{\Tr}{Tr}
\DeclareMathOperator{\rank}{rank}
\newcommand{\R}{\mathbb R}
\newcommand{\C}{\mathbb C}
\newcommand{\Z}{\mathbb Z}
\newcommand{\N}{\mathbb N}
\newcommand{\Q}{\mathbb Q}
\begin{document}
\section*{022\quad 有限Galois扩张的正规基与群代数模消去}
\subsection*{来源与收录范围}
J. S. Milne, \emph{Fields and Galois Theory}, v5.00, June 2021。主命题 Definition 5.17、Theorem 5.18，印刷第66页；只选其 \emph{Uniform proof}，印刷第67--68页，及第66页群代数模的等价表述。此处 Galois 扩张按 Definition 5.17 的有限范围使用。取得日期：2026-09-28。
\par\url{https://www.jmilne.org/math/CourseNotes/FT500.pdf}
\par\url{https://www.jmilne.org/math/CourseNotes/ft.html}
\subsection*{作者命题与人类证明}
\begin{definition}[5.17]
Let $E$ be a finite Galois extension of $F$. A basis for $E$ as an $F$-vector space is called a normal basis if it consists of the conjugates of a single element of $E$. In other words, a normal basis is one of the form
\[
\{\sigma\alpha\mid\sigma\in\operatorname{Gal}(E/F)\}
\]
for some $\alpha\in E$.
\end{definition}
\begin{theorem}[5.18, Normal basis theorem]
Every Galois extension has a normal basis.
\end{theorem}
The group algebra $FG$ of a group $G$ is the $F$-vector space with basis the elements of $G$ endowed with the multiplication extending that of $G$. Every $F$-linear action of $G$ on an $F$-vector space $V$ extends uniquely to an action of $FG$ on $V$.

Let $E/F$ be a Galois extension with Galois group $G$. Then $E$ is an $FG$-module, and Theorem 5.18 says that there exists an element $\alpha\in E$ such that the map
\[
\sum_{\sigma}a_\sigma\sigma\longmapsto\sum_{\sigma}a_\sigma\sigma\alpha:FG\longrightarrow E
\]
is an isomorphism of $FG$-modules, i.e., that $E$ is a free $FG$-module of rank $1$.

\begin{proof}[Uniform proof]
Recall that a module is indecomposable if it is nonzero and cannot be written as a direct sum of two nonzero submodules. The Krull--Schmidt theorem says that every nonzero module $M$ of finite length over a ring can be written as a direct sum of indecomposable modules and that the indecomposable modules occurring in a decomposition are unique up to order and isomorphism. Thus $M=\bigoplus_i m_iM_i$ where $M_i$ is indecomposable and $m_iM_i$ denotes the direct sum of $m_i$ copies of $M_i$; the set of isomorphism classes of the $M_i$ is uniquely determined and, when we choose the $M_i$ to be pairwise nonisomorphic, each $m_i$ is uniquely determined. From this it follows that two modules $M$ and $M'$ of finite length over a ring are isomorphic if $mM\simeq mM'$ for some $m\geq1$.

Consider the $F$-vector space $E\otimes_F E$. We let $E$ act on the first factor, and $G$ act on the second factor (so $a(x\otimes y)=ax\otimes y$, $a\in E$, and $\sigma(x\otimes y)=x\otimes\sigma y$, $\sigma\in G$). We'll prove Theorem 5.18 by showing that
\[
\underbrace{FG\oplus\cdots\oplus FG}_{n}\simeq E\otimes_F E
\simeq\underbrace{E\oplus\cdots\oplus E}_{n}
\]
as $FG$-modules ($n=[E:F]$).

For $\sigma\in G$, let $\lambda_\sigma:E\otimes_F E\to E$ denote the map $x\otimes y\mapsto x\cdot\sigma y$. Then $\lambda_\sigma$ is obviously $E$-linear, and $\lambda_\sigma(\tau z)=\lambda_{\sigma\tau}(z)$ for all $\tau\in G$ and $z\in E\otimes_F E$. I claim that $\{\lambda_\sigma\mid\sigma\in G\}$ is an $E$-basis for $\Hom_{E\text{-linear}}(E\otimes_F E,E)$. As this space has dimension $n$, it suffices to show that the set is linearly independent. But if $\sum_\sigma c_\sigma\lambda_\sigma=0$, $c_\sigma\in E$, then
\[
0=\sum_\sigma c_\sigma\bigl(\lambda_\sigma(1\otimes y)\bigr)
=\sum_\sigma c_\sigma\cdot\sigma y
\]
for all $y\in E$, which implies that all $c_\sigma=0$ by Dedekind's theorem 5.14.

Consider the map
\[
\Phi:E\otimes_F E\longrightarrow EG,\qquad
z\longmapsto\sum_\sigma\lambda_\sigma(z)\cdot\sigma^{-1}.
\]
Then $\Phi$ is $E$-linear. If $\Phi(z)=0$, then $\lambda_\sigma(z)=0$ for all $\sigma\in G$, and so $z=0$ in $E\otimes_F E$ (because the $\lambda_\sigma$ span the dual space). Therefore $\Phi$ is injective, and as $E\otimes_F E$ and $EG$ both have dimension $n$ over $E$, it is an isomorphism. For $\tau\in G$,
\[
\Phi(\tau z)=\sum_\sigma\lambda_\sigma(\tau z)\cdot\sigma^{-1}
=\sum_\sigma\lambda_{\sigma\tau}(z)\cdot\tau(\sigma\tau)^{-1}
=\tau\Phi(z);
\]
and so $\Phi$ is an isomorphism of $EG$-modules. Thus
\[
E\otimes_K E\simeq EG\simeq FG\oplus\cdots\oplus FG
\]
as an $FG$-module.

On the other hand, for any basis $\{e_1,\ldots,e_n\}$ for $E$ as an $F$-vector space,
\[
E\otimes_F E=(e_1\otimes E)\oplus\cdots\oplus(e_n\otimes E)
\simeq E\oplus\cdots\oplus E
\]
as $FG$-modules. This completes the proof.
\end{proof}

\subsection*{整理说明（非作者正文）}
本题只计正规基定理一个目标，不把有限域、无限域及统一证明分别计数。本题允许 Krull--Schmidt 有限长度模分解与消去、Dedekind 特征标线性无关性及有限 Galois 扩张的基本性质作为标准先修；实际承担的是两个模分解、$E$-线性对偶基以及群作用相容的具体同构构造。难度依据这些结构之间的连接，不依据定理名称或篇幅。

为使转录中的两个不同符号可区分，把作者的同构映射统一排作 $\Phi$，另一个群元素排作 $\tau$，不改变映射或运算。原文结论一行写作 $E\otimes_K E$，前后底域均为 $F$；这里照录 $K$ 并在此指出该处原文符号不一致，不把修订伪装成作者原文。保留作者统一证明的顺序，没有加入新的证明。已取得 PDF 文字；第68页截图接口失败，未声称逐页图像核对，原 PDF 未保存到本地。
\subsection*{可修改的简短标注（非作者正文）}
$c$：有限 Galois 扩张的正规基

$a$：群代数模；标量扩张与张量积；Dedekind 特征标线性无关性；对偶基；Krull--Schmidt 分解与消去

\texttt{cross\_theory\_status = yes}。把共轭元素构成基的问题转成 $E\simeq FG$ 的模同构；经 $E\otimes_F E$ 的两个分解与等变映射取得多重同构，再用有限长度模消去返回正规基。位置：Uniform proof，印刷第67--68页。

全部标注均为非穷尽、可修改初稿；未标注不表示未使用。
\subsection*{许可与修改记录}
Copyright 1996--2021 J. S. Milne. 原书 v5.00 的第2页明确以 Creative Commons Attribution--NonCommercial--ShareAlike 4.0 许可发布。本文摘录及附加整理沿用该许可，不暗示原作者认可选题或标签。2026-09-28：助手转录为独立 TeX，并加入来源、范围与标注。许可：\url{https://creativecommons.org/licenses/by-nc-sa/4.0/}。
\end{document}
